%--------------------------------------------------------------
\section{Comparison with earlier work}\label{sec:compare}
%--------------------------------------------------------------
\subsection{Hodge theory (formal coefficients)}
Gualtieri proved a generalized Hodge star, the Born--Infeld inner product, Laplacian equalities, the Hodge decomposition and the $dd^{\Jc}$-lemma on generalized K\"ahler manifolds (\tC\ Gualtieri, \emph{Generalized geometry and the Hodge decomposition}, arXiv:math/0409093; Gualtieri, \emph{Generalized K\"ahler geometry}, Comm.\ Math.\ Phys.\ 331 (2014) 297--331, arXiv:1007.3485; Cavalcanti, \emph{New aspects of the $dd^c$-lemma}, arXiv:math/0501406, and J.\ Geom.\ Phys.\ 57 (2006) 121--132. Theorem numbers are omitted because these were not read closely \tQ). These are theories with trivial or $H$-twisted coefficients; a Bochner--Kodaira--Nakano identity with line-bundle (curved) coefficients has not been written, as far as surveyed. The results of this note can be positioned as filling this gap for $\beta$-transforms, not by Hodge theory but by comparing logarithmic complexes with constructible sheaves.

\subsection{Pal--Poddar}
\tC\ Pal--Poddar, \emph{Strong generalized holomorphic principal bundles}, arXiv:2404.18113, treat compact generalized complex manifolds of regular type (constant type $k$) whose symplectic leaf space is a K\"ahler orbifold, and give Hodge theory and Serre duality for strong generalized holomorphic (SGH) bundles (transversal Dolbeault isomorphism Thm.~8.9, Cor.~8.10 (numbers \tQ); Hodge decomposition and Serre duality Thm.~11.2, 11.4, 11.6, 11.7), K\"ahler identities (Prop.~11.8, Lemma~11.9; numbers \tQ), and Kodaira-type and Serre-type vanishing theorems for positive SGH line bundles (Def.~11.10, 11.11; Thm.~11.12). The hypotheses of regular type and orbifold leaf space are used to guarantee that the foliation is a smooth leaf bundle and to guarantee transversal Hodge theory (ellipticity). Differences from the present note:
\begin{itemize}
\item The type-jumping case (the type jumps on the degeneracy divisor of $\sigma$) is outside the framework of Pal--Poddar; examples such as $\mathbb{CP}^2$ in this note are of this kind.
\item SGH bundles form a proper subclass of generalized holomorphic bundles (flat along the leaves and holomorphic transversally). The Poisson modules of this note are general GH bundles on a $\beta$-transformed structure.
\item The statement and proof structure of Thm.~11.12 were checked by re-reading \S11: $H^q(M,(\mathcal G^*M)^p\otimes E)=0$ for $p+q>k$, and $H^q(M,E'\otimes E^m)=0$ for $m\ge m_0$, $q>0$. The hypotheses (compactness, a K\"ahler orbifold leaf space, positivity of $E$) are used for transversal Hodge theory and for positivity of the curvature term of a Bochner--Kodaira--Nakano type identity.
\end{itemize}
Doubts about the original text noticed in this survey (likely misprints or omissions in the source, but unverified \tQ): (i) the statement carries a factor $(i/2\pi)$ while the proof has no $1/2\pi$ (harmless); (ii) the sign discrepancy about $gc_1$ comes from the differing definitions of the characteristic polynomial in (9.9) and of $\det(I\pm tA/2\pi i)$ in \S10.1; (iii) ``a positive class gives a metric of positive curvature'' does not follow from Def.~11.11 and is assumed in the proof of Thm.~11.12; (iv) Def.~11.11 does not state compactness or the K\"ahler condition; (v) the proof of Lemma~11.9(1) is a computation in pointwise local frames; (vi) $T^{1,0}$ versus $T^{1,0*}$ in Thm.~13.4(2)(3), and the binomial coefficient for $q=n$ in Example 13.10.

\subsection{The work of Tengzhou Hu}
\tC\ T.\ Hu, \emph{Kodaira vanishing theorems for K\"ahler Lie algebroids}, arXiv:2403.10876;\\ T.\ Hu, \emph{Lie groupoid Riemann--Roch--Hirzebruch theorem and applications}, arXiv:2403.13424.
\begin{itemize}
\item \emph{Setting} (2403.13424): a Lie groupoid $G\rightrightarrows M$ (Def.~1; hypotheses of properness or $s$-connectedness are not stated) and its Lie algebroid $\mathfrak g$. A ``complex Lie algebroid'' (Def.~2) is a \emph{real} Lie algebroid $\mathfrak g$ with an integrable complex structure $J$ in the fibre direction (this is different from the complex Lie algebroid $L$ as a subbundle of $\bT_\C$). Kähler Lie algebroids are Def.~11, positive bundles Def.~12.
\item \emph{Claims} (2403.13424): Thm.~4 (Riemann--Roch) states, for $M$ compact, $\dbar_{\mathfrak g}$ elliptic (claimed, but no proof was found) and $\Omega$ a $\mathfrak g$-invariant density (unimodularity),
\[\mathrm{Ind}_\Omega(\dbar_{\mathfrak g}+\dbar_{\mathfrak g}^*)=(2\pi i)^{-k}\int_M(-1)^{\mathrm{rk}\,\mathfrak g}\bigl\langle\ch^{\mathfrak g}(\mathfrak g^{*p,0})\ch^{\mathfrak g}(E)\Td^{\mathfrak g}_\C(\mathfrak g),\Omega_{\mathfrak g}\bigr\rangle.\]
Thm.~7 (vanishing): if $E$ is positive and $p+q>n$, the \emph{kernel} of the $\mathfrak g$-Dolbeault Laplacian vanishes (vanishing of cohomology does not follow: for a general Lie algebroid $\mathfrak g$ the Laplacian is not necessarily elliptic and Hodge theory is unavailable). In 2403.10876, a K\"ahler Lie algebroid is Def.~2, the K\"ahler identities are Thm.~1 and Prop.~4, the BKN identity is Lemma~4, and the main vanishing theorem is in the introduction.
\item \emph{The generalized Dolbeault complex is outside the scope}\ \tP\ (derived; comparison with the source only to the extent indicated below): $L$ is a complex Lie algebroid without $J$, and $\wedge L^*$ in general has no $(p,q)$ bigrading. For a holomorphic Poisson structure $\sigma$ the real Lie algebroid $A_{\mathbb R}$ falls within Hu's setting, but its Dolbeault complex is $\wedge\mathfrak g^{*0,q}$ built from $\bar\sigma$, which differs from the complex $(T^{0,1}\bowtie A^{1,0})$ of this note.
\item \emph{Does the index see $m$?}\ \tK: Hu's index is a groupoid index, and is expected to depend polynomially on $m$, roughly like $\int\ch(E)\Td(T^*)\ch(\wedge^pT)$, whereas the total Euler characteristic of this note is independent of $m$ (Prop.~\ref{prop:euler}). Moreover the residue $-m/3$, which governs the periodic jumps, is not part of Hu's data. Hence Hu's theorem is not ``Riemann--Roch for generalized Dolbeault cohomology''.
\item \emph{Inconsistencies in the source}: 2403.13424 (obtained through summaries): $n$ is undefined in Thm.~7; the sign constants differ among Thm.~5, Prop.~1 and Thm.~6; ``unimodular'' is defined in two ways. 2403.10876 (checked by reading further and part of the appendix \tC): Lemma~5 gives $[L,\Lambda]=(p+q-n)\,\mathrm{id}$, but the proof of Thm.~2 writes $\Delta_{\bar\partial_{\mathfrak g}}=\Delta^{1,0}-[L,\Lambda]=\Delta^{1,0}+(p+q-n)\,\mathrm{id}$, and the sign in the last equality does not match ($-(p+q-n)=n-p-q$). Moreover the substitution $\nabla^2=\sqrt{-1}L$ means $\sqrt{-1}R=-L$, opposite to the definition of ``positive'' ($\sqrt{-1}R>0$, Def.~13). The statement itself (the kernel vanishes for positive $E$ when $p+q>n$) agrees with classical Kodaira--Nakano; so the printed proof has a sign slip, not the claim. Misprints in Cor.~1 (arrows and the symbols $s/t$). The appendix (Lemmas 6--10, the computation of $[\bar\partial_{\mathfrak g},\iota_\omega]$) could not be verified because the indices are hard to read. The reference list and publication data are unchecked. \tQ
\item \emph{Correction to the earlier version of this note (preliminary survey)}: the earlier version gave the author as ``S.~Hu'' and a different title; both were wrong, and the above is correct.
\end{itemize}

\subsection{Integrability and groupoids}
\tC\ Gualtieri--Li (arXiv:1206.3674, IMRN 2014) classify Hausdorff symplectic groupoids of log-symplectic manifolds and show that the Lie algebroid $T_X(-\log D)$ of logarithmic vector fields tangent to a smooth $D$ is integrable. Bailey--Gualtieri (arXiv:1611.03850) show the equivalence between integration of a generalized complex structure to a weak holomorphic symplectic groupoid and integrability of the real Poisson structure (Thm.~7.1: integrability of the underlying real Poisson structure is necessary and sufficient; Thm.~7.1 itself says nothing about Hausdorffness or $s$-connectedness, and the construction Thm.~7.6 uses an $s$-connected, $s$-simply connected integration). Laurent-Gengoux--Sti\'enon--Xu (arXiv:0803.2031): Thm.~3.17, 3.22. Crainic--Fernandes (arXiv:math/0105033): Thm.~4.1, Cor.~5.2. These provide groupoids as input for index theory, but do not treat integration of the complex Lie algebroid $L$ itself.

%--------------------------------------------------------------
\section{Open problems and next steps}\label{sec:open}
%--------------------------------------------------------------
\begin{enumerate}
\item \textbf{The treasure not yet touched}: a generalization of the Kodaira embedding theorem. Since $\HH^0$ consists of Casimir sections (Observation~\ref{obs:casimir}), it must be formulated not as an embedding of $M$ itself but as a map to the quotient separated by Casimir sections (e.g.\ the leaf space). The next concrete step is to study base points and separation properties of the linear system of Casimir sections in the $\mathbb{CP}^2$ example ($\HH^0=\C F^{m/3}$, i.e.\ 1-dimensional for $m\in3\Z_{\ge0}$).
\item \textbf{Generalization of Serre-type vanishing} (Conjecture~\ref{conj:A3}): a definition of a ``positive'' generalized holomorphic line bundle, and a construction of a filtration playing the role of the $E_1$ term for $\Jc$ that are not $\beta$-transforms. The comparison with Pal--Poddar Thm.~11.12 was done by re-reading \S11 (the groups differ; the transversal-degree reading is consistent).
\item \textbf{Formulation of Riemann--Roch}: instead of the total Euler characteristic, one must fix a refinement (the $E_1$ term of a filtration by the type of $\Jc$, a periodic quasi-polynomial, the indicial roots $H^*(\partial U)$). Proving Conjecture~\ref{conj:I3} requires an index theory for the complex Lie algebroid $L$ (unimodularity, and a derivation of Fredholmness $\Leftrightarrow r\notin\Z$). An APS-type extension to generalized complex manifolds with boundary is a still more distant problem.
\item \textbf{Beyond log-symplectic}: proofs for the local theory of non-nodal singularities (only quasi-homogeneous germs were computed), a global formula for unequal residues, non-reduced $D$ (infinite-dimensional stalks and linear growth), and non-toric examples in dimension $\ge4$ (see Observation~\ref{obs:arith}).
\item \textbf{Generalized K\"ahler side}: write down the Bochner--Kodaira--Nakano identity with line-bundle coefficients in the generalized K\"ahler framework. The interaction of the $H$-twist with the curvature term may be an obstacle \tK.
\item \textbf{Resolution of the unverified items of this note}: the \tQ\ items of \S\ref{sec:ledger}. In particular Mokri's theorem, a re-derivation of LGSX Lemma~4.22, a reading of Gualtieri 2011 \S4, and the remaining bibliographic data.
\end{enumerate}

\paragraph{Priorities for the next steps (author's proposal).}
(a) As a comparison case in which the type does not jump, compute the same ``Poincar\'e polynomial'' for the regular-type examples of Pal--Poddar and compare with the periodic jumps of this note. (b) For the $\mathbb{CP}^2$ example, use $\HH^0=\C F^{m/3}$ to study the geometry of the linear system of Casimir sections and look for the ``right form'' of a Kodaira-type theorem. (c) As a basis for an index-theoretic proof of Conjecture~\ref{conj:I3}, formulate rigorously the relation between the weighted $b$-complex on the real Lie algebroid $T_X(-\log D)$ and the residue $r$.

